% Header include for markdown-to-Tufte-PDF via pandoc + xelatex.
% Loaded with --include-in-header, so it lands at the END of the preamble
% (after the tufte-handout class and any \setmainfont pandoc emits). Keep
% font selection OUT of here — it belongs in frontmatter / the wrapper so a
% user-supplied mainfont is not clobbered.

% --- Sidenotes -------------------------------------------------------------
% The whole point of this skill: route pandoc's footnotes into the margin.
% tufte-handout keeps \footnote as a bottom-of-page note and offers a separate
% \sidenote. Pandoc emits \footnote{...} for markdown footnotes, so alias it
% to \sidenote globally. A note-style CSL (e.g. chicago-note-bibliography)
% then also lands its citations in the margin, which is the Tufte ideal.
%
% Caveat: \sidenote uses \marginpar, which LaTeX forbids inside floats and
% some boxes. A footnote placed inside a table cell or a figure caption will
% therefore error. See SKILL.md ("Sidenotes and the float caveat").
\let\footnote\sidenote

% --- Code blocks -----------------------------------------------------------
% Pandoc defines a "Highlighting" Verbatim environment for syntax-highlighted
% code. Redefine it via fvextra so long lines wrap instead of overflowing the
% (narrow, by Tufte design) text column.
\usepackage{fvextra}
\DefineVerbatimEnvironment{Highlighting}{Verbatim}{%
  breaklines,%
  breakanywhere,%
  fontsize=\footnotesize,%
  commandchars=\\\{\}%
}

% --- Graphics --------------------------------------------------------------
% tufte-handout loads graphicx, but load it unconditionally so the raw-LaTeX
% figure/marginfigure pattern documented in SKILL.md works even when the doc
% contains no markdown images.
\usepackage{graphicx}

% --- Tables ----------------------------------------------------------------
% Pandoc auto-loads longtable + booktabs for markdown tables; load tabularx
% and booktabs unconditionally so hand-written raw-LaTeX tables work too.
% NOTE: longtable does not coexist with the Tufte margin layout — long tables
% should go in a \begin{fullwidth}...\end{fullwidth} block. See SKILL.md.
\usepackage{tabularx}
\usepackage{booktabs}
